\section{Study Intervention}

This study compares two randomized arms under a single integrated protocol. Arm A,
the experimental arm, combines an investigational drug, perioperative daraxonrasib
(RMC-6236) administered at the recommended Phase 2 dose (RP2D), with an
investigational device, the on-premises large language model (LLM) directed
eight-arm robotic pancreaticoduodenectomy (Whipple) system (PancreSpeed II). Arm B,
the active control, is modified FOLFIRINOX \cite{Conroy2018} plus
institutional-standard high-volume pancreaticoduodenectomy that is not
LLM-directed. This section describes both arms, their administration, their
preparation and accountability, the measures taken to minimize bias, the compliance
machinery for each, and the permitted and prohibited concomitant therapy. In Arm A
the drug and device are coupled at one decision point only: the perioperative
pause-and-restart advisory for the drug is informed by the operative and pathology
results of the device procedure (\S6.1.2).

\subsection{Study Intervention Administration}

\subsubsection{Study Intervention Description}

\paragraph{Experimental arm drug: daraxonrasib (RMC-6236).}
Daraxonrasib is an orally bioavailable, RAS(ON) multi-selective (pan-RAS)
small-molecule inhibitor that binds the active, guanosine triphosphate bound state
of mutant RAS in complex with cyclophilin A, suppressing downstream effector
signaling across the common KRAS oncoprotein variants found in PDAC. The agent
received FDA Breakthrough Therapy designation in June 2025 for previously treated
metastatic PDAC \cite{revbreakthrough}, and its dose is anchored to the pan-KRAS
RASolute 302 program, in which 300~mg once daily is the established RP2D
\cite{rasolute302,kawchak2026phase1}. In Arm A daraxonrasib is administered orally,
once daily, in a perioperative schedule: a defined pre-operative course, a
protocol-mandated pause spanning the operative window, and a postoperative restart
whose timing is governed by the advisory described in \S6.1.2. The drug is
investigational in this perioperative, curative-intent surgical context and for the
drug-device use here.

\paragraph{Control arm: modified FOLFIRINOX plus standard pancreaticoduodenectomy.}
Arm B delivers institutional-standard care. The systemic component is modified
FOLFIRINOX, the multi-agent regimen established in the resectable and adjuvant PDAC
setting \cite{Conroy2018}, administered on its standard schedule with the
institution's supportive-care practice. The surgical component is a high-volume
pancreaticoduodenectomy performed by the institutional standard, robotic or open,
that is not LLM-directed and does not use the eight-arm investigational platform.
Arm B thus represents the contemporary standard of care against which the
experimental drug-device regimen is compared.

\paragraph{Experimental device: on-premises LLM-directed eight-arm Whipple (PancreSpeed II).}
The investigational device is an eight-arm parallel coelomic surgical platform in
which an on-premises repository-pinned LLM acts as a second-opinion advisory oracle
held strictly outside the robot vendor kinematic stack, never as an autonomous
controller. The platform carries 56 mechanical degrees of freedom ($8\times7$
per-arm degrees of freedom) and a 640-channel mixed-rate sensor stack (80 channels
per arm). Force channels sample at 100 kHz; all remaining channels sample at the
10 kHz command rate. Positioning accuracy is 0.05 mm root mean square at a peak tip
velocity of $1{,}200$ mm/s. The hard cyber-physical safety envelope comprises a
per-arm tip-force cap of $\leq 3$ N, a cumulative cross-arm tip-force cap of
$\leq 18$ N enforced at the heartbeat boundary, a 10 kHz heartbeat coordination bus
broadcasting a 64-byte frame on a $100\,\mu$s per-arm watchdog deadline, an
emergency arm park within $50\,\mu$s of a missed or out-of-parameter frame, and a
cross-arm emergency stop (E-stop) that halts the full platform in $\leq 3$ ms; a
software-independent hardware E-stop backs the software chain and acts within
$\leq 500$ ms. A five-vessel no-fly gate enforces standoff distances around the
major peripancreatic vessels (1.0 mm for the superior mesenteric vein and portal
vein, 1.5 mm for the hepatic artery, celiac axis, and superior mesenteric artery, a
3.0 mm soft margin, and a 5.0 mm hard-stop margin). The platform is multicenter
fleet-harmonized so that every instance behaves identically across the eight
centers. In this Phase 2 protocol the device operates at Stage 2 supervised
semiautonomy under continuous human oversight with immediate manual override; full
autonomy is prohibited consistent with 21 CFR \S312.21(e). The advisory control
loop and platform architecture are shown in Figure 11 and Figure 12.

\begin{mermaidfig}
\node[mmin]  (sens) at (0,0)      {640 sensor\\channels\\100 kHz force /\\10 kHz other};
\node[mmstep](map)  at (4.2,0)    {Sensor to x,y,z\\Cartesian command\\frame};
\node[mmgoal](llm)  at (8.4,0)    {On-premises LLM\\advisory, hash-pinned\\to commit + seed};
\node[mmdec] (gate) at (8.4,-2.8) {Safety gate\\vessel zones + force\\caps; 10 kHz beat,\\$100\,\mu$s watchdog};
\node[mmin]  (vend) at (4.2,-2.8) {Robot vendor\\kinematic stack\\(isolated from LLM)};
\node[mmgoal](act)  at (0,-2.8)   {8-arm actuators\\$\leq 3$ N/arm,\\$\leq 18$ N cumulative};
\node[mmdark](aud)  at (0,-5.2)   {Hash-chained\\federated audit\\21 CFR part 11};
\draw[mmedge] (sens) -- (map);
\draw[mmedge] (map) -- (llm);
\draw[mmedge] (llm) -- (gate);
\draw[mmedge] (gate) -- (vend);
\draw[mmedge] (vend) -- (act);
\draw[mmedge] (act.south) -- (aud.north);
\draw[mmedge] (act.west) -- ++(-0.5,0) |- ($(sens.west)+(-0.5,0)$) -- (sens.west);
\draw[mmedged] (gate.south) -- ++(0,-0.15) -| (act.south);
\end{mermaidfig}
\vspace{-0.7cm}
\figcaption{Figure 11. On-premises LLM advisory control loop. Sensor data is mapped
to a Cartesian command\\frame; the LLM issues hash-pinned advisory commands from
outside the vendor kinematic stack\\(second-opinion oracle isolation); and a safety
gate enforces the vessel zones, the force caps, the 10 kHz\\ heartbeat, and the
$100\,\mu$s watchdog before any actuator moves. Actuator telemetry returns to the
sensor stack, and every step is written to a hash-chained federated audit trail
under 21 CFR part 11 \cite{ecfr2026part11}.}

\begin{mermaidfig}
\node[mmgoal](hub) at (4.2,0)     {10 kHz heartbeat\\coordination bus\\64-byte frame,\\$100\,\mu$s deadline};
\node[mmstep](a1)  at (0,-2.4)    {Arm 1 / Arm 2\\hybrid u-w-p\\dissection (P1-P8)};
\node[mmstep](a3)  at (4.2,-2.4)  {Arm 3 bipolar +\\retractor; Arm 4\\NIR + ICG (P1-P8)};
\node[mmstep](a5)  at (8.4,-2.4)  {Arm 5 anastomosis\\(P5-P7); Arm 6 coag\\(P2, P5-P8)};
\node[mmstep](a7)  at (8.4,-5)    {Arm 7 suction +\\irrigation (P1-P8)};
\node[mmstep](a8)  at (4.2,-5)    {Arm 8 imaging +\\drain (P1, P4, P8)};
\node[mmin]  (spec)at (0,-5)      {56 DOF ($8\times7$)\\640 channels (80/arm)\\0.05 mm RMS at\\$1{,}200$ mm/s};
\draw[mmedge] (hub) -- (a1);
\draw[mmedge] (hub) -- (a3);
\draw[mmedge] (hub) -- (a5);
\draw[mmedge] (a5) -- (a7);
\draw[mmedge] (a3) -- (a8);
\draw[mmedged] (a1) -- (spec);
\end{mermaidfig}
\vspace{-0.7cm}
\figcaption{Figure 12. Eight-arm platform architecture. Eight cooperating arms
(56 degrees of freedom, 640 sensor\\channels at 80 per arm, 0.05 mm root mean square
positioning at $1{,}200$ mm/s) are coordinated by a\\10 kHz heartbeat bus; each
arm's tool assignment and operative-phase coverage drive the Schedule of
Activities.}

The per-arm tool assignment crossed with the eight operative phases (P1 through P8)
is given in Table~\ref{tab:arms}, and ten representative channels drawn from the
640-channel inventory are given in Table~\ref{tab:sensors}.

\begin{table}[htbp]
\centering
\footnotesize
\caption{Per-arm tool assignment and operative-phase coverage for the eight-arm LLM-directed Whipple platform (PancreSpeed II).}
\label{tab:arms}
\begin{tabularx}{\textwidth}{L{1.2cm} L{4.9cm} L{4.0cm} Y}
\toprule
\textbf{Arm} & \textbf{Primary tool} & \textbf{Phase coverage} & \textbf{Sensor channels (10 of 80)} \\
\midrule
Arm 1 & Hybrid u-w-p end-effector & P1 to P8 (active dissect) & Tip force, joint torque, ID \\
Arm 2 & Hybrid u-w-p end-effector & P1 to P8 (active dissect) & Tip force, joint torque, ID \\
Arm 3 & Bipolar + retractor & P1 to P8 (vessel control) & Tip force, current, force \\
Arm 4 & NIR + ICG imaging probe & P1 to P8 (imaging) & ICG fluorescence, NIR depth \\
Arm 5 & Anastomosis probe & P5, P6, P7 (idle others) & Ring tension, ring strain \\
Arm 6 & Bipolar coag + cautery & P2, P5, P6, P7, P8 & Tip force, RF power \\
Arm 7 & Suction + irrigation & P1 to P8 & Suction pressure, pH \\
Arm 8 & Final images + drain placement & P1, P4, P8 (light load) & ICG fluorescence, drain ID \\
\bottomrule
\end{tabularx}
\end{table}

\begin{table}[htbp]
\centering
\footnotesize
\caption{Ten representative sensor channels drawn from the 640-channel inventory (80 channels per arm).}
\label{tab:sensors}
\begin{tabularx}{\textwidth}{L{1.0cm} L{3.8cm} L{3cm} L{2.4cm} Y}
\toprule
\textbf{Arm} & \textbf{Channel name} & \textbf{Rate} & \textbf{Unit} & \textbf{Dynamic range} \\
\midrule
1 & tip\_force\_x & 100 kHz & N & 0 to 3.0 N \\
1 & joint\_torque\_3 & 10 kHz & N$\cdot$m & $-50$ to $50$ \\
2 & tip\_force\_z & 100 kHz & N & 0 to 3.0 N \\
3 & retractor\_force & 10 kHz & N & 0 to 5.0 N \\
4 & icg\_fluorescence & 10 kHz & a.u. & 0 to $65{,}535$ \\
4 & nir\_depth & 10 kHz & mm & 0 to 50 \\
5 & ring\_tension & 10 kHz & N & 0 to 1.0 N \\
6 & rf\_power & 10 kHz & W & 0 to 80 \\
7 & suction\_pressure & 10 kHz & kPa & $-90$ to 0 \\
8 & drain\_id & 10 kHz & enum & 0 to 4 \\
\bottomrule
\end{tabularx}
\end{table}

\paragraph{Physical AI System Description fields (21 CFR \S312.23(g)).}
Because the device participates in the investigational drug delivery pathway, the
combined IND/IDE carries a full Physical AI System Description meeting the eleven
fields of 21 CFR \S312.23(g) \cite{kawchak2026adapt312}. The fields, and the
location of the corresponding evidence in this submission, are summarized below;
the multicenter conduct adds fleet harmonization and a federated audit across the
eight centers.

\begin{itemize}[leftmargin=1.4em,itemsep=1pt,topsep=2pt]
\item[(i)] \textit{System identification and classification:} manufacturer, model
designation, software version, surgical robot category, and the current Unified
Safety Level (USL) score computed across simulation switching, AI integration,
cross-robot sharing, and clinical-trial collaboration.
\item[(ii)] \textit{Role in drug administration:} secondary role (the device does
not itself dispense daraxonrasib; it conducts the resection and reconstruction and
generates the operative and pathology inputs to the restart advisory), with
permitted autonomy levels and human-oversight requirements for each procedure.
\item[(iii)] \textit{Hardware specifications:} 56 degrees of freedom, 0.05 mm root
mean square positioning at $1{,}200$ mm/s, 100 kHz force sensing, the $\leq 3$ N
per-arm and $\leq 18$ N cumulative tip-force limits, and end-effector details.
\item[(iv)] \textit{Software and AI/ML specifications:} the operating and
middleware stack, the control algorithms, and the on-premises advisory LLM together
with its training data characteristics, validation methodology, performance
metrics, and Predetermined Change Control Plan (PCCP).
\item[(v)] \textit{Simulation validation data:} the upgraded Phase 0 campaign
results, including at least 5000 simulated procedures reproduced across at least
three independent frameworks and the tightened sim-to-real gap quantification
(trajectory gap $< 1.5$ mm, tip-force gap $< 0.4$ N).
\item[(vi)] \textit{Digital twin specifications:} the digital-twin methodology,
calibration sources, update frequency, and validation metrics used for operative
rehearsal and intra-operative guidance.
\item[(vii)] \textit{Safety systems:} the emergency-stop mechanisms (hardware and
software), the force and torque limits, the five-vessel no-fly gate, collision
detection, workspace-boundary enforcement, fail-safe behaviors, and the
pre-procedure safety matrix.
\item[(viii)] \textit{Cybersecurity assessment:} the network architecture,
deny-by-default authentication and authorization, encryption, vulnerability
management, and incident response.
\item[(ix)] \textit{MCP integration:} the Model Context Protocol server topology,
conformance level, robot capability profile, and the hash-chained federated
audit-trail implementation under 21 CFR part 11 \cite{ecfr2026part11}.
\item[(x)] \textit{Human-robot interaction protocols:} the operator control
stations, status displays, alarm systems, handoff procedures, and the training
requirements for operating and supervising personnel across the fleet.
\item[(xi)] \textit{Maintenance and calibration:} the scheduled maintenance
intervals, calibration procedures, out-of-service criteria, per-site qualification,
and re-qualification requirements (\S6.2).
\end{itemize}

\subsubsection{Dosing and Administration}

\paragraph{Daraxonrasib dosing.}
In Arm A daraxonrasib is administered orally once daily at the fixed RP2D of 300~mg.
Because the RP2D is already established in the Phase 1 predicate and anchored to the
RASolute 302 pan-KRAS program \cite{kawchak2026phase1,rasolute302}, there is no 3+3
dose escalation in this Phase 2 study; every participant in Arm A receives the same
fixed daily dose, and the trial tests comparative efficacy rather than dose-finding.

\paragraph{Perioperative pause-and-restart advisory.}
In Arm A daraxonrasib is paused before surgery and restarted postoperatively on a
day chosen by an advisory that is LLM-bound and human-confirmed, never autonomous.
The advisory takes two inputs: the pancreaticojejunostomy ISGPS postoperative
pancreatic-fistula grade (A, B, or C) \cite{Bassi2017} and the serum daraxonrasib
trough relative to a 0.5 ng/mL threshold. It recommends a restart on postoperative
day T+7, T+14, or T+21 (or a continued hold): T+7 for a Grade A fistula with a
sub-threshold trough, T+14 for a Grade B fistula or a delayed serum rise, and T+21
or continued hold for a Grade C fistula, which also maps to the drug stop rule of
\S7. The advisory logic is shown in Figure 13.

\begin{mermaidfig}
\node[mmin]  (pre) at (5.2,0)     {Pre-operative hold\\daraxonrasib paused\\before surgery};
\node[mmstep](in1) at (0,-2.4)    {Input: PJ ISGPS\\fistula grade\\(A / B / C)};
\node[mmstep](in2) at (10.4,-2.4) {Input: serum trough\\vs 0.5 ng/mL};
\node[mmdec] (adv) at (5.2,-2.4)  {LLM-bound\\restart advisory\\(human confirmed)};
\node[mmgoal](r7)  at (0,-5.0)    {Restart T+7\\Grade A,\\trough $< 0.5$};
\node[mmstep](r14) at (5.2,-5.0)  {Restart T+14\\Grade B / delayed\\serum rise};
\node[mmdark](r21) at (10.4,-5.0) {Restart T+21 /\\hold; Grade C};
\draw[mmedge] (pre) -- (in1);
\draw[mmedge] (pre) -- (in2);
\draw[mmedge] (in1) -- (adv);
\draw[mmedge] (in2) -- (adv);
\draw[mmedge] (adv) -- (r7);
\draw[mmedge] (adv) -- (r14);
\draw[mmedge] (adv) -- (r21);
\end{mermaidfig}
\vspace{-0.7cm}
\figcaption{Figure 13. Daraxonrasib perioperative pause-and-restart advisory in
Arm A. The restart day is keyed to the pancreaticojejunostomy ISGPS fistula grade
and the serum trough relative to 0.5 ng/mL: T+7 for Grade A with a sub-threshold
trough, T+14 for Grade B or a delayed serum rise, and T+21 or continued hold for
Grade C. The advisory is bound and human-checked, not autonomous.}

\paragraph{Device administration: operative phases and the three anastomoses.}
The device has no pharmacologic dose; its administration is the conduct of the
pancreaticoduodenectomy across eight operative phases (P1 through P8) under
continuous human oversight. Phases P1 through P4 comprise the dissection and
specimen removal (Arms 1 and 2 dissecting under Arm 3 vessel control and Arm 4
near-infrared indocyanine-green imaging); the specimen is removed at the end of P4
and reconstruction begins. The three reconstructive anastomoses are then constructed
in sequence under closed-loop ring-tension control by the Arm 5 anastomosis probe,
which issues a soft warning whenever measured ring tension falls outside the
controller band: the duct-to-mucosa pancreaticojejunostomy (PJ) in P5 with a band of
0.30 to 0.60 N, the end-to-side hepaticojejunostomy (HJ) in P6 with a band of 0.20
to 0.50 N, and the antecolic gastrojejunostomy (GJ) in P7 with a band of 0.40 to
0.80 N; P8 covers final imaging, hemostasis, and drain placement. The
pancreaticojejunostomy is the dominant determinant of postoperative
pancreatic-fistula risk and is therefore the anastomosis on which the fistula-grade
input to the drug advisory is based. The reconstruction sequence and the
ring-tension bands are shown in Figure 14.

\begin{mermaidfig}
\node[mmin]  (rem) at (0,0)       {Specimen removed\\(end P4)\\reconstruction begins};
\node[mmgoal](pj)  at (4.2,0)     {Pancreatico-\\jejunostomy (PJ)\\duct-to-mucosa, P5\\0.30 to 0.60 N};
\node[mmgoal](hj)  at (8.4,0)     {Hepatico-\\jejunostomy (HJ)\\end-to-side, P6\\0.20 to 0.50 N};
\node[mmgoal](gj)  at (8.4,-2.8)  {Gastro-\\jejunostomy (GJ)\\antecolic, P7\\0.40 to 0.80 N};
\node[mmstep](mon) at (4.2,-2.8)  {Ring-tension monitor\\(Arm 5)\\soft-warn outside\\band};
\draw[mmedge] (rem) -- (pj);
\draw[mmedge] (pj) -- (hj);
\draw[mmedge] (hj) -- (gj);
\draw[mmedged] (mon) -- (pj);
\draw[mmedged] (mon) -- (hj);
\draw[mmedged] (mon.east) -- ++(0.4,0) |- (gj.west);
\end{mermaidfig}
\vspace{-0.7cm}
\figcaption{Figure 14. Three reconstructive anastomoses and their controller
ring-tension bands: the duct-to-mucosa pancreaticojejunostomy (0.30 to 0.60 N), the
end-to-side hepaticojejunostomy (0.20 to 0.50 N), and the antecolic
gastrojejunostomy (0.40 to 0.80 N), each guarded by the Arm 5 ring-tension monitor
issuing a soft warning outside the band.}

\subsection{Preparation, Handling, Storage, and Accountability}

\subsubsection{Acquisition and Accountability}

Daraxonrasib is supplied by the sponsor to the investigational-product pharmacy at
each participating site and is dispensed only by the site pharmacist or an
authorized designee against a written prescription tied to an enrolled,
Arm A participant. The pharmacy maintains a drug-accountability log recording, for
every unit, the lot number, the quantity received, the quantity dispensed, the
quantity returned, and the quantity destroyed, reconciled against shipment records
at each monitoring visit. Returned and unused product is retained for reconciliation
and destroyed only after sponsor authorization. The modified FOLFIRINOX components
of Arm B are obtained and accounted for through standard institutional oncology
pharmacy practice. The investigational device is provided under the IDE and is not a
consumable; its accountability is established through the asset record, the
per-procedure pre-procedure safety matrix, and the hash-chained federated audit
trail that captures every operative use across the fleet (\S6.4). Custody of all
components is restricted to authorized, trained personnel, and all transfers are
documented.

\subsubsection{Formulation, Appearance, Packaging, and Labeling}

Daraxonrasib is provided as oral tablets in the strengths required to compose the
300~mg once-daily RP2D. Tablets are packaged in child-resistant, light-protective
containers labeled in accordance with 21 CFR \S312.6, bearing the caution statement
``Caution: New Drug. Limited by Federal (or United States) law to investigational
use,'' together with the protocol number, lot number, storage conditions, and a
unique container identifier. Because this is an open-label study, the
participant-facing label identifies the product and the assigned daily dose; no
label information is used to influence the masked outcome assessors described in
\S6.5. The investigational device and its single-use accessories are labeled per
21 CFR part 812 with the investigational-device caution statement and are not
relabeled or repackaged at the site.

\subsubsection{Product Storage and Stability}

Daraxonrasib is stored at controlled room temperature in a secured,
access-controlled, temperature-monitored location within the investigational-product
pharmacy, in its original light-protective packaging, until dispensed. Continuous
temperature monitoring with documented excursion handling is maintained; any
excursion outside the labeled range is quarantined and is not dispensed until the
sponsor confirms continued suitability against the stability data on file. Stability
is supported by the sponsor's stability program, and product is dispensed only
within its labeled expiry or retest date. The investigational device is stored and
maintained in its qualified clinical environment between procedures under the
maintenance and calibration program below, with the same qualification applied
identically at every site in the harmonized fleet.

\subsubsection{Preparation and Device Maintenance and Calibration}

\paragraph{Drug preparation.}
No compounding or reconstitution is required for daraxonrasib; preparation consists
of the pharmacist verifying the participant identity, the RP2D, and the expiry, then
dispensing the once-daily oral dose with administration instructions. Tablets are
swallowed whole with water and are not crushed, split, or chewed unless a future
protocol amendment provides validated instructions. Modified FOLFIRINOX is prepared
per institutional oncology pharmacy standards.

\paragraph{Device maintenance and calibration (21 CFR \S312.23(g)(xi)).}
Before each clinical procedure the system must pass the pre-procedure safety matrix:
verification of the 640-channel sensor stack, confirmation of the per-arm force
calibration against the $\leq 3$ N per-arm and $\leq 18$ N cumulative caps, exercise
of the 10 kHz heartbeat with its $100\,\mu$s watchdog, confirmation of the
five-vessel no-fly gate, a successful software and hardware E-stop self-test, and
confirmation that the positioning accuracy remains within the 0.05 mm root mean
square specification at $1{,}200$ mm/s. Scheduled maintenance, periodic force and
positional calibration, per-site qualification, and fleet harmonization verification
are performed at the manufacturer-specified intervals and documented. A system that
fails any pre-procedure check, exhibits an out-of-tolerance calibration, or
experiences a hard-stop E-stop event is taken out of service and is returned to
clinical use only after re-qualification against the same matrix and sponsor
authorization. All maintenance, calibration, and out-of-service events are recorded
in the hash-chained federated audit trail.

\subsection{Measures to Minimize Bias: Randomization and Blinding}

Allocation uses central permuted-block randomization 1:1, stratified by resectability
(resectable versus borderline-resectable), KRAS allele (G12D versus other G12),
neoadjuvant therapy (yes versus no), and site, so that the central mechanism prevents
allocation bias and balances the strata across arms \cite{Schulz2010consort}.
Treatment is open-label because a surgical-system intervention cannot be masked from
the operating team or the participant; bias control is therefore concentrated
downstream of the intervention. The imaging-based endpoints pass through blinded
independent central review (BICR); the R0 resection status and the major pathologic
response (MPR) are determined on central masked pathology by reviewers who are masked
to the assigned arm, to the intra-operative device telemetry, and to the LLM advisory
record, with discordant or indeterminate margins adjudicated by an independent masked
second pathologist; and all confirmatory endpoints pass through blinded endpoint
adjudication independent of the operative team. In this way the open-label delivery
does not propagate into outcome ascertainment.

\subsection{Study Intervention Compliance}

\paragraph{Daraxonrasib adherence.}
In Arm A, adherence is documented through the dispensing and accountability log,
returned-tablet counts at each visit, and a participant dosing diary reconciled
against the pharmacy record. Adherence is corroborated objectively by the serum
daraxonrasib assay, which both confirms exposure and supplies the trough value used
by the restart advisory (\S6.1.2); a documented trough discordant with the reported
intake prompts adherence counseling and is recorded in the case report form.
Adherence to modified FOLFIRINOX in Arm B is captured through standard cycle
administration records. Material non-adherence in either arm is captured as a
protocol deviation and informs the per-protocol analysis (\S9).

\paragraph{Device compliance.}
Device compliance is enforced before and during every procedure rather than
reconstructed afterward. No operative phase may begin until the pre-procedure safety
matrix passes in full, and the 10 kHz heartbeat, the $100\,\mu$s watchdog, the force
caps, and the five-vessel no-fly gate operate continuously throughout the case. Every
advisory command, every safety-gate verdict, every force and trajectory sample at the
recorded rates, and every E-stop event is written to a 21 CFR part 11 hash-chained,
append-only federated audit trail \cite{ecfr2026part11}, so that compliance with the
cyber-physical envelope is verifiable for each case across the fleet; any deviation is
detectable and time-stamped. The audit trail also satisfies the preservation
requirement that applies around a reportable Physical AI adverse event (\S8).

\subsection{Concomitant Therapy}

All medications, biologics, blood products, radiotherapy, and clinically significant
non-pharmacologic interventions taken from the time of informed consent through the
end of safety follow-up are recorded in the case report form with indication, dose,
route, and dates. Standard perioperative supportive care is permitted and encouraged
in both arms, including antiemetics, analgesia, venous thromboembolism prophylaxis,
peri-operative antimicrobial prophylaxis, proton-pump inhibitors, and
pancreatic-enzyme replacement as clinically indicated, with attention in Arm A to
agents that may interact with daraxonrasib exposure. Prohibited during study
participation are systemic anticancer therapy outside the assigned regimen (for
Arm A, cytotoxic chemotherapy, other targeted agents, or immunotherapy beyond the
protocol; for Arm B, agents beyond the modified FOLFIRINOX standard), other
investigational agents or investigational devices, and any concomitant use of a
second surgical robotic system for the index operation. Strong inducers or inhibitors
of the daraxonrasib principal metabolic pathway are avoided where feasible in Arm A;
when an interacting agent is clinically unavoidable, its use is documented and
reviewed against exposure data.

\subsubsection{Rescue Medicine}

Rescue is available along two axes. For surgical and device-related complications,
the operative team may at any time convert to manual or open technique (\S7) and
deliver standard rescue care: transfusion and hemostatic resuscitation for bleeding,
vasoactive support for hemodynamic instability, percutaneous or operative drainage and
broad-spectrum antimicrobials for a clinically relevant postoperative pancreatic
fistula or intra-abdominal collection, and somatostatin-analogue therapy where
indicated for high-output fistula. For drug toxicity, daraxonrasib in Arm A is held or
discontinued per the stopping rules (\S7) and managed symptomatically with the
appropriate supportive agents (for example antidiarrheals and aggressive dermatologic
care for the rash and gastrointestinal toxicities characteristic of the class), and
modified FOLFIRINOX in Arm B is dose-modified or held per its established toxicity
management, with dose interruption and reduction governed by the protocol. All rescue
interventions are recorded as concomitant therapy and, where the criteria are met, as
adverse or serious adverse events.
